\chapter{CUDA 内存、地址空间与指针}
\label{app:cuda-memories-address-spaces-pointers}

\begin{center}
  \small Isaac Gelado 和 Mark Harris 特别贡献
\end{center}

本书大部分内容都假设异构计算系统中的主机与设备采用一种相当简单的交互模型。如第 2 章
所述，在这一简单模型中，每台设备都有自己的设备内存（CUDA 全局内存），与主机内存（即
系统内存）彼此分离。设备上运行的内核要处理的数据，必须调用 \code{cudaMemcpy()} 函数从
主机内存传输到设备内存。设备生成的数据也必须先调用 \code{cudaMemcpy()} 函数从设备内存
传回主机内存，主机才能使用。早期 GPU 世代只支持这种简单的主机/设备交互模型。该模型虽然
简单易懂，却在应用层面带来了几个问题。

第一，磁盘控制器和网络接口卡等 I/O 设备经过专门设计，能够高效操作主机内存。由于设备
内存与主机内存分离，输入数据必须从主机内存传到设备内存，输出数据则必须从设备内存传到
主机内存，I/O 设备才能使用。这些额外传输会增加 I/O 延迟，降低可达到的 I/O 吞吐量。对
许多应用而言，如果 I/O 设备能直接操作设备内存，既可提高应用总体性能，也可简化应用代码。

第二，Python 等传统编程系统把应用数据结构放在主机内存中，其中有些数据结构非常大。与
主机内存相比，支持 CUDA 的 GPU 设备内存往往带宽更高、容量却更小，迫使应用开发者把大型
数据结构划分成若干块，使每块都能装入设备内存。例如第 21 章把三维静电势能网格数组划分为
二维切片，在主机内存和设备内存之间传输。对许多应用而言，如果整个数据结构都能驻留在设备
内存中，效果会好得多。某些应用的数据结构甚至没有合适的分块方法；对于这些应用，最好让
GPU 能直接访问主机内存中的数据，或者让 CUDA 运行时系统软件迁移内核执行期间使用的数据。

随着越来越多的应用采用 GPU 计算，这些需求推动 CUDA 系统软件开发者和 GPU 硬件设计者
提供更好的解决方案。研究人员从 CUDA 早期就已意识到这些需求并提出解决方案
\cite{appc-gelado}。本附录余下各节将简要回顾为克服上述限制而取得的进展。

\section{统一设备内存地址空间}
\label{sec:appc-unified-device-address-space}

在早期支持 CUDA 的 GPU 中，共享内存、局部内存和全局内存各自形成独立的地址空间。开发者
可以使用指向全局内存的指针，却不能对其他内存使用指针。从 2009 年推出的 Fermi 架构开始，
这些内存成为统一地址空间的一部分。借助这一统一地址空间，同一组加载/存储指令和指针地址
就能访问任意 GPU 内存空间（全局、局部或共享内存），无须为每一种内存使用不同的指令和
指针。这样更容易抽象掉某个操作数究竟位于哪种内存的问题，程序员只需在分配时处理这一
区别；无论 CUDA 数据对象来自哪一内存区域，把它们传给其他过程和函数也变得更简单。

统一设备内存地址空间支持使 CUDA 代码模块更具“可组合性”。也就是说，CUDA 设备函数可以
接收一个可能指向上述任意一种内存的指针。例如，如果没有统一 GPU 地址空间，那么对某个
实参可能驻留的每种内存类型，设备函数都需要分别实现一个版本。统一 GPU 地址空间让所有
主要 GPU 内存类型中的变量都能以相同方式访问，因此一个设备函数可以接收驻留于不同 GPU
内存类型的实参。若函数实参指针指向共享内存位置，代码运行得较快；若指向全局内存位置，
代码运行得较慢。程序员仍可把手工放置和传输数据作为一种性能优化。这项能力显著降低了构建
生产级 CUDA 库的成本，也让 CUDA 内核函数和设备函数获得了完整的 C 和 C++ 指针支持。

\section{对主机内存的零拷贝访问}
\label{sec:appc-zero-copy}

2009 年，CUDA 2.2 引入了对主机内存的零拷贝（zero-copy）访问。主机代码由此可以向内核
提供一个特殊的、指向主机内存的设备数据指针。GPU 上运行的 CUDA 内核代码可以使用该指针，
经由 PCIe 总线等系统互连直接访问主机内存，而无须调用 \code{cudaMemcpy()}。零拷贝内存
属于固定页主机内存（见第 23 章），通过调用 \code{cudaHostAlloc()} 分配，并把标志实参设为
\code{cudaHostAllocMapped}。前文曾提到，该标志实参的其他取值用于更高级的用法。
\code{cudaHostAlloc()} 返回的数据指针不能直接传给内核；主机代码必须先调用
\code{cudaHostGetDevicePointer()} 取得有效的设备数据指针，再把该函数返回的设备数据指针
传给内核。这意味着，主机代码和设备代码使用不同的数据指针访问主机内存中的同一份数据。

正如第 23 章所述，为供 GPU 零拷贝访问而分配的主机内存页必须固定下来，以免 GPU 访问期间
操作系统意外地把数据换出。显然，这种访问会受到系统互连长延迟和有限带宽的影响。系统互连
带宽通常不到全局内存带宽的 10\%。正如第 5 章所学，除非采用分块技术大幅减少每次浮点操作
对应的全局内存访问次数，否则内核性能通常受全局内存带宽限制。如果内核的大部分内存访问
都面向零拷贝内存，系统互连带宽对内核执行速度的限制会更加严重。因此，只应把零拷贝内存
用于 GPU 内核偶尔且稀疏访问的应用数据结构。

\section{统一虚拟地址空间}
\label{sec:appc-uvas}

2011 年，CUDA 4 引入了源自 GMAC 库的统一虚拟地址空间（Unified Virtual Address Space，
UVAS）\cite{appc-gelado}。在这一版 CUDA 之前，主机和设备各有自己的虚拟地址空间，分别把
主机或设备数据指针映射到主机或设备的物理内存位置。这些彼此分离的虚拟地址空间意味着，
同一个物理内存位置可能在主机和设备上通过不同的虚拟地址访问；使用零拷贝内存时就是如此。
UVAS 规定主机和设备共享单一虚拟地址空间，并保证每个物理内存地址只映射到一个虚拟内存
位置。这样，CUDA 运行时只需检查数据指针的虚拟内存地址，就能判断它引用的是主机内存还是
设备内存。凭借这一特性，调用 \code{cudaMemcpy()} 时不再需要指定数据复制方向；程序员只需
指定 \code{cudaMemcpyDefault}，系统便会根据指针值推断传输方向。

需要特别指出，CUDA 4 中的 UVAS 并不保证指针所引用的数据一定可以访问。例如，主机代码
不能使用 \code{cudaMalloc()} 返回的设备指针直接访问设备内存，反之亦然。对主机内存的零
拷贝访问是一个例外：主机代码可以把指向零拷贝内存（通过 \code{cudaHostAlloc()} 分配）的
指针作为内核参数传给设备。内核代码解引用这一零拷贝指针时，指针值会被转换成物理系统内存
位置，并用于经 PCIe 总线直接访问系统内存。注意，这种方法不一定允许内核代码解引用从某个
内存位置读出的指针值，例如在遍历链式数据结构时沿指针链前进；只有当所有指针都引用通过
\code{cudaHostAlloc()} 分配的内存时，才能这样做。

零拷贝内存在可支持的数据结构类型和数据访问带宽两方面都有局限，因此需要在 UVAS 之外
继续改进 GPU 架构的内存模型。

\section{大型 GPU 虚拟地址空间和物理地址空间}
\label{sec:appc-large-address-spaces}

早期支持 CUDA 的 GPU 有一项根本限制，即虚拟地址和物理地址的大小。这些早期设备支持 32
位虚拟地址以及最高 32 位物理地址。因此，设备内存大小被限制在 4 GB，也就是 32 个物理
地址位能够寻址的最大内存量。此外，无论数据集驻留在主机内存还是设备内存，CUDA 内核都只能
处理小于 4 GB 的数据集，因为 32 位指针最多只能访问这么多虚拟内存位置。与此同时，现代
CPU 使用 64 位虚拟地址，但实际使用其中 48 位。GPU 使用的 32 位虚拟地址无法容纳这些主机
虚拟地址，这也限制了零拷贝内存访问所能支持的数据结构类型，并使这种访问必须采用
\code{cudaHostAlloc()} 内存分配机制。

为消除这一限制，从 2013 年推出的 Kepler GPU 架构开始，各代 GPU 采用了现代虚拟内存
架构，使用 64 位虚拟地址以及至少 40 位的物理地址。显而易见的好处是，这些 GPU 可以配备
超过 4 GB 的 DRAM，CUDA 内核也可以处理大型数据集。扩大的虚拟地址空间和物理地址空间
不仅显然支持更大的设备内存，也为更好的主机/设备交互模型打开了大门。例如，无论一份数据
位于主机内存还是设备内存，主机和设备现在都能用完全相同的指针值访问它。

\section{统一物理地址空间}
\label{sec:appc-unified-physical-address-space}

大型 GPU 物理地址空间允许 CUDA 系统软件把系统中不同 GPU 的设备内存放入统一物理地址
空间。其好处是，一个 GPU 只需解引用映射到另一个 GPU 物理地址的数据指针，就能直接访问
连接在同一 PCIe 总线和/或 NVLink 域中的任意其他 GPU 的内存。在 Kepler GPU 架构之前，
不同 GPU 之间的通信（例如第 23 章模板示例中的 halo 交换）只能通过主机代码触发的设备到设备
内存复制完成。这样既要消耗额外内存存放从其他 GPU 复制来的数据，又会因主机与设备同步以及
设置内存复制所需的操作产生额外性能开销。能够直接访问系统中其他设备的内存后，只需把设备
指针作为内核调用参数传入，内核代码就能用它加载和/或存储驻留在另一 GPU 设备内存中的数据。

\section{统一内存}
\label{sec:appc-unified-memory}

2013 年，CUDA 6 引入了统一内存（Unified Memory）。它创建由 CPU 和 GPU 共享的托管
虚拟内存页池，弥合 CPU 与 GPU 之间的鸿沟。CPU 和 GPU 可以通过同一个指针访问托管内存。
托管内存中的变量可以驻留在 CPU 物理内存、GPU 物理内存中，甚至同时驻留在两者中。CUDA
运行时软件和硬件负责实现数据迁移和一致性支持。最终效果是：在 CPU 上运行的代码看来，
托管内存就像 CPU 内存；在 GPU 上运行的代码看来，它又像 GPU 内存。当然，应用必须执行
屏障或原子操作等适当的同步操作，协调对托管内存位置的任何并发访问。共享的全局虚拟地址
空间使应用中的所有变量都能从 CPU 或 GPU 代码以相同的虚拟地址访问。编程工具和运行时
系统把这种内存架构呈现给应用后，可以带来若干重大好处。

\begin{figure}[htbp]
  \centering
  \begin{minipage}[t]{0.47\textwidth}
    \centering\textbf{CPU 代码}
\begin{lstlisting}[language=C++,basicstyle=\ttfamily\scriptsize]
void sortfile(FILE *fp, int N) {
    char *data;
    data = (char *)malloc(N);
    fread(data, 1, N, fp);
    qsort_char(data, N, 1);

    use_data(data);
    free(data);
}
\end{lstlisting}
  \end{minipage}\hfill
  \begin{minipage}[t]{0.49\textwidth}
    \centering\textbf{采用统一内存的 CUDA 6 代码}
\begin{lstlisting}[language=C++,basicstyle=\ttfamily\scriptsize]
void sortfile(FILE *fp, int N) {
    char *data;
    cudaMallocManaged(&data, N);
    fread(data, 1, N, fp);
    qsort_char<<<...>>>(data, N, 1);
    cudaDeviceSynchronize();
    use_data(data);
    cudaFree(data);
}
\end{lstlisting}
  \end{minipage}
  \caption{统一内存简化了把 CPU 代码（左）移植为 CUDA 代码（右）的过程。依据原书图 C.1
  重排。}
  \label{fig:c-1}
\end{figure}

统一内存的一项好处，是减少把 CPU 代码移植到 CUDA 所需的工作量。图 \ref{fig:c-1} 左侧
给出一个简单的 CPU 代码示例。借助统一内存，只需作两项简单改动即可把它移植到 CUDA。
第一项改动是用 \code{cudaMallocManaged()} 和 \code{cudaFree()} 取代 \code{malloc()} 和
\code{free()}。第二项改动是不再调用 \code{qsort()} 函数，而是启动内核并执行设备同步。
显然，程序员仍需编写或者取得一个并行 \code{qsort} 内核。这里要说明的是，主机代码的改动
直接明了，也易于维护。

CUDA 6 统一内存的性能受 Kepler 和 Maxwell GPU 架构的硬件能力限制。每次启动网格前，
必须把 CPU 修改过的所有托管内存位置的内容刷新到 GPU 设备内存。CPU 和 GPU 不能同时访问
同一块托管内存分配，统一内存地址空间也受 GPU 物理内存大小限制。出现这些限制，是因为这些
GPU 架构无法支持主机内存与设备内存之间的一致性，而且数据迁移大多由软件完成。

\section{访问完整的主机虚拟地址空间}
\label{sec:appc-full-host-virtual-address-space}

2016 年，Pascal GPU 架构把 GPU 地址转换能力扩展到可以容纳 49 位虚拟地址。这个范围足以
覆盖现代 CPU 的 48 位虚拟地址空间以及 GPU 自身的内存。统一内存程序由此可以把系统中所有
CPU 和 GPU 的完整地址空间作为单一虚拟地址空间访问，不再受 GPU 地址转换机制所能处理的
虚拟地址位数限制。于是，CPU 和 GPU 可以真正共享指针值，GPU 也能遍历主机内存中基于指针
的数据结构。

\section{缺页处理}
\label{sec:appc-page-fault-handling}

从 Pascal GPU 架构开始提供的内存缺页处理支持，是使统一内存功能更加无缝的一项关键特性。
缺页处理能力与全系统虚拟地址空间结合，消除了 CUDA 系统软件在每次网格启动之前同步（刷新）
全部托管内存内容的需要。主机或设备修改托管内存中的数据时，可以使对方的数据副本失效；
CUDA 运行时由此能够实现一致性机制。副本失效可通过页面映射和保护机制完成。启动网格时，
CUDA 系统软件不再需要把托管内存数据在 GPU 中的全部副本更新到最新状态。如果网格访问的
一份数据在设备内存中的副本已被主机标为失效，GPU 就会处理一次缺页，把数据更新到最新状态，
然后恢复执行。

如果 GPU 上运行的线程网格访问一个未驻留在其设备内存中的页面，也会发生缺页，从而按需自动
把该页面迁移到 GPU 内存。另一种办法是，如果预计只会偶尔访问数据，就把页面映射进 GPU
地址空间，经系统互连访问；按访问进行映射有时会比迁移更快。这项支持把零拷贝访问从主机
内存中的有限部分推广到整个主机内存乃至系统中其他 GPU 的内存。注意，统一内存作用于整个
系统：GPU（以及 CPU）可以因缺页而从 CPU 内存或系统中其他 GPU 的内存迁移内存页。如果
CPU 函数解引用一个指针，并访问映射到 GPU 物理内存的变量，这一数据访问仍会得到服务，只是
延迟可能较长。凭借这种能力，CUDA 程序可以更容易地调用尚未移植到 GPU 的遗留库。在此前的
CUDA 内存架构中，开发者若要使用遗留库函数在 CPU 上处理数据，必须手工把数据从设备内存
传输到主机内存。

具有缺页处理能力的统一内存提供了一种比原始零拷贝访问通用得多的 CPU/\allowbreak GPU
交互机制。它
允许 GPU 遍历主机内存中的大型数据结构。从 Pascal 架构开始，即使链式数据结构并不驻留在
零拷贝内存中，GPU 设备也能遍历它。这是因为主机代码和设备代码使用相同的指针值引用同一个
变量。因此，设备可以沿主机构建的链式数据结构中嵌入的指针值进行遍历，反之亦然。在 CAD
等某些应用领域，主机物理内存系统可能拥有数百 GB 的容量。这些物理内存系统不可或缺，因为
应用要求整个数据集都“驻留于内存”。GPU 能够直接访问极大的 CPU 物理内存以后，加速这些
应用就成为可能。

\section{虚拟地址空间控制}
\label{sec:appc-virtual-address-space-control}

CUDA 11 引入了一组底层 API，使程序员可以更灵活地分配内存。新 API 允许使用
\code{cuMemAddressReserve()} 预留一段虚拟地址空间。之后，程序员可以使用
\code{cuMemCreate()} 在任意设备上分配物理内存，再使用 \code{cuMemMap()} 把它映射到
预留范围中的任意位置。借助这些 API，可以跨多台设备构建自定义数据结构布局。例如，可以
把一个三维体数据分配到多台设备上，同时用单一指针引用它。

\begin{chapterbibliography}{1}
  \bibitem{appc-gelado}
  I. Gelado, J. E. Stone, J. Cabezas, S. Patel, N. Navarro, W.-m. W. Hwu,
  An asymmetric distributed shared memory model for heterogeneous parallel systems,
  in: \emph{Proceedings of the Fifteenth International Conference on Architectural
  Support for Programming Languages and Operating Systems}, 2010, pp.~347--358.
\end{chapterbibliography}
